\documentclass[10pt]{article}
\usepackage[margin=0.62in]{geometry}
\usepackage[T1]{fontenc}
\usepackage{lmodern}
\usepackage{microtype}
\usepackage{xcolor}
\usepackage{multicol}
\usepackage{enumitem}
\usepackage{booktabs}
\usepackage{fancyhdr}
\usepackage{mdframed}
\usepackage{titlesec}
\usepackage[hidelinks]{hyperref}

\definecolor{ink}{HTML}{17324D}
\definecolor{signal}{HTML}{D8613C}
\definecolor{wash}{HTML}{F2F5F7}
\definecolor{muted}{HTML}{5D6B75}
\pagestyle{fancy}
\fancyhf{}
\lhead{\sffamily\small\color{ink}\textbf{FIELD NOTES}}
\rhead{\sffamily\small\color{muted}VOL. 02 · ISSUE 04}
\cfoot{\sffamily\small\color{muted}\thepage}
\renewcommand{\headrulewidth}{0.4pt}
\renewcommand{\headrule}{\hbox to\headwidth{\color{signal}\leaders\hrule height \headrulewidth\hfill}}
\titleformat{\section}{\sffamily\large\bfseries\color{ink}}{}{0pt}{}[\color{signal}\titlerule]
\setlist{nosep,leftmargin=*}
\newmdenv[linecolor=signal,backgroundcolor=wash,linewidth=1pt,skipabove=8pt,skipbelow=8pt,innerleftmargin=9pt,innerrightmargin=9pt]{dispatch}
\newcommand{\dateline}[1]{\small\sffamily\color{muted}#1\par\medskip}
\newcommand{\pullquote}[1]{\begin{dispatch}\sffamily\large\bfseries\color{ink}#1\end{dispatch}}

\begin{document}
\begin{center}
  {\sffamily\fontsize{30}{32}\selectfont\bfseries\color{ink} FIELD NOTES}\par
  \smallskip{\sffamily\color{signal}\large A compact newsletter starter for teams that learn in public}\par
  \medskip{\sffamily\small\color{muted} April 2026 \quad · \quad Three columns, one clear point of view}
\end{center}
\vspace{-4pt}\hrule\vspace{10pt}

\begin{multicols*}{3}
\section*{The short version}
\dateline{EDITORIAL / 04 APRIL}
Small experiments compound when a team records what changed, what surprised it, and what happens next. This issue is a reusable pattern for sharing progress without turning every update into a meeting.

\pullquote{Make the next useful action visible.}

\section*{A better weekly rhythm}
The most useful project updates are neither status reports nor polished announcements. They are small, dated observations that let another person pick up the thread.

\begin{enumerate}
  \item Name the question before the result.
  \item Show one piece of evidence.
  \item Record the decision and its owner.
\end{enumerate}

\section*{Three signals to watch}
\begin{tabular}{@{}llr@{}}
\toprule
Signal & This week & Change\\
\midrule
Open questions & 12 & $-3$\\
Field interviews & 18 & $+5$\\
Decisions logged & 7 & $+2$\\
\bottomrule
\end{tabular}

\columnbreak
\section*{In focus}
\dateline{A SMALLER FEEDBACK LOOP}
The neighbourhood repair clinic used to collect suggestions on a single form. Volunteers now sort each suggestion into a question, a proposed trial, or a request for help. The categories are simple enough to use at the counter and specific enough to make the next conversation concrete.

After four weeks, the team had fewer duplicate requests and a clearer handoff between the people greeting visitors and the people scheduling work. The change was not a new system; it was a shared vocabulary.

\subsection*{What made it stick}
\begin{itemize}
  \item A two-minute review at the end of each shift.
  \item One visible owner for every open question.
  \item A Friday note that closed the loop with contributors.
\end{itemize}

\section*{A useful question}
\begin{dispatch}
\sffamily\textbf{Where does information disappear between observation and action?}\\[4pt]
\small Identify the handoff, then change only that handoff for one week.
\end{dispatch}

\columnbreak
\section*{Around the desk}
\dateline{TOOLS / PEOPLE / PRACTICE}
\textbf{Read}\\
Write the one-page version before writing the long version. The missing detail will tell you what to investigate.

\medskip\textbf{Try}\\
End every meeting with a sentence beginning “By Friday, we will know whether...”.

\medskip\textbf{Share}\\
Bring one unresolved question to the next newsletter instead of waiting for a perfect answer.

\section*{Calendar}
\begin{itemize}
  \item 08 Apr — prototype review
  \item 11 Apr — field-note deadline
  \item 15 Apr — community demo
  \item 18 Apr — monthly retro
\end{itemize}

\section*{One last note}
This layout is intentionally modular. Replace sections with announcements, profiles, event listings, or short research summaries. Keep the masthead and the three-column rhythm; the content can change every issue.

\vfill
\begin{center}\sffamily\small\color{muted}\href{https://example.org}{example.org/field-notes}\end{center}
\end{multicols*}
\end{document}
